\documentclass[11pt,a4paper]{article}

\usepackage[margin=2.5cm]{geometry}
\usepackage{amsmath}
\usepackage{booktabs}
\usepackage{fancyhdr}
\usepackage{xcolor}
\usepackage[hidelinks]{hyperref}
\usepackage{enumitem}
\usepackage{titlesec}

\pagestyle{fancy}
\fancyhf{}
\lhead{\small Brian Mwai -- 1050555}
\rhead{\small B.Sc. Physics -- Proposal Amendment}
\cfoot{\thepage}
\renewcommand{\headrulewidth}{0.4pt}

\titleformat{\section}{\large\bfseries}{\thesection}{0.8em}{}
\titleformat{\subsection}{\normalsize\bfseries}{\thesubsection}{0.8em}{}
\setlength{\parskip}{0.5em}
\setlength{\parindent}{0pt}

\newcommand{\naa}{\mathrm{NAA}}

\begin{document}

\begin{center}
{\large THE CATHOLIC UNIVERSITY OF EASTERN AFRICA}\\[0.2em]
{\normalsize Faculty of Science \quad$\cdot$\quad Department of Physics}\\[1.4em]
{\Large\bfseries Proposal Amendment}\\[0.6em]
{\large \emph{Measuring the Invisible:} A Sociophysics Framework for Quantifying\\
Emotionally Manipulative Media Content Using Predictive Neural Encoding}\\[1.4em]
\end{center}

\begin{center}
\begin{tabular}{ll}
\textbf{Student} & Brian Mwai (1050555) \\
\textbf{Programme} & Bachelor of Science in Physics \\
\textbf{Supervisor} & Dr. Songa Mutambi \\
\textbf{Instrument} & \url{https://monarch.brianmwai.com} \\
\textbf{Status} & Requires supervisor sign-off before the thesis is written \\
\end{tabular}
\end{center}

\vspace{1em}
\hrule
\vspace{1em}

\section{What this document is}

The April proposal was written before the model at the centre of it had been run.
Once it was run, three things turned out to be different from what the proposal
assumed, and one dataset the proposal names is no longer available to anybody.

None of these is a mistake in the work. Each is a property of the tools, found by
testing them. This document states each one in plain terms, says what the project
does instead, and asks for four decisions.

\textbf{The short version:} the project still delivers the instrument the proposal
promised, and answers three of its four research questions. What it can no longer
do is measure the amygdala, because the released model does not predict it.

\section{The four changes}

\subsection{Change 1: the model cannot see the amygdala}

\textbf{What the proposal assumed.} That TRIBE v2 predicts activity across 29{,}286
brain targets including the amygdala and nucleus accumbens. Equation (4) defines the
project's central measurement as the ratio of amygdala activity to prefrontal activity.

\textbf{What the released model actually does.} It predicts 20{,}484 points on the
outer surface of the brain only. The amygdala sits deep inside the brain, so it is
not among them. This is not a setting that can be switched on; those outputs do not
exist in the released weights. It is stated in the checkpoint's own configuration
file, which is included with the model.

\textbf{Why this matters in one sentence.} The quantity Equation (4) is built from
cannot be computed at all.

\textbf{What the project does instead.} It measures the same balance one level up,
on the brain surface: regions associated with emotional salience versus regions
associated with deliberate reasoning. Everywhere this appears, it is called a
\emph{cortical proxy} and never called the amygdala.

\textbf{Honest consequence.} The thesis can say that emotionally charged content is
predicted to shift the balance of surface brain activity. It cannot say anything
about the amygdala itself. That limitation is stated in the abstract, not buried.

\subsection{Change 2: the original formula does not work on real content}

\textbf{What the proposal specified.} A ratio,
\[
\naa \;=\; \frac{A_{\text{amyg}}}{A_{\text{PFC}} + \delta},
\qquad \delta = 10^{-3}.
\]

\textbf{What happens when it is used.} The model reports activity in standardised
units, centred on zero. That means the bottom of the fraction is negative roughly
half the time. Dividing by a negative number flips the answer's sign; dividing by
a number near zero makes it explode. On a 40-item test run the formula produced no
usable value for a single item.

\textbf{What the project does instead.} It subtracts rather than divides:
\[
\naa_{\text{signed}} \;=\; A_{\text{aff}} - A_{\text{del}}.
\]
Positive means the emotional side leads. Negative means the reasoning side leads.
This is always defined. The cost is that the number is no longer dimensionless: it
carries the units of the model's standardised output, and the thesis says so.

\textbf{Acknowledgement.} This is the change Dr. Mutambi raised at the proposal
defence. It is recorded here as her suggestion, adopted once running the model
showed why the ratio could not survive standardised output.

\subsection{Change 3: the calibration found nothing, twice}

\textbf{What the proposal specified.} Objective (iv): fit the coupling constant
$\hat{\alpha}$, the number connecting the brain index to the physics model's field
strength.

\textbf{What happened.} It was fitted twice, on independent runs. Both times the
confidence interval comfortably included zero, and the fit performed worse than
guessing on held-out data. In plain terms: there is no evidence at this sample size
that the connection exists at a size this study could measure.

\textbf{What the project does instead.} It reports the null result and quotes no
value for $\hat{\alpha}$ anywhere -- not in a table, not in a caption, not ``for
illustration''. The phase structure is then presented across a range of possible
$\alpha$ values, which shows how the physics would behave for each, and states that
the data do not pin it down.

\textbf{Why this is acceptable physics.} A null result that is properly powered
constrains what everyone working after you has to believe. It is only worthless when
the sample is too small to have detected anything, which is why the sample size and
the detectable effect size are both reported.

\subsection{Change 4: a named dataset was withdrawn by its authors}

\textbf{What the proposal specified.} NELA-GT-2021, for the credibility scores in
objective (iv).

\textbf{What happened.} Its authors removed it from public access on 1 January 2024,
after the proposal was written. It now requires per-application manual approval.
This was verified rather than assumed: with a valid Harvard Dataverse account, a
control file from an unrelated dataset downloads normally while every NELA file
returns ``forbidden'', and the dataset does not appear in search at all.

\textbf{What the project does instead.} It uses SemEval-2019 Task 4 (Hyperpartisan
News Detection), openly licensed CC BY 4.0. This is arguably a better fit: NELA
labels whether a \emph{publisher} is credible, while this labels the \emph{article},
and the project's categories are defined by what an article does.

\section{What the thesis will deliver}

\subsection{Delivered}

\begin{enumerate}[leftmargin=1.6em]
\item \textbf{The measurement instrument, working and public.} The proposal calls
this the primary deliverable and says so three times. It exists: the batch pipeline,
the four-category corpus, and the corpus-level audit report described in
\S5.8(iii) -- ranked table of every item, distribution per category, free-energy
landscape, and per-item flags. Released as open-source code with the corpus manifest
and one complete example report. Live at \url{https://monarch.brianmwai.com}.

\item \textbf{The physics, derived in full.} The mean-field Ising treatment, the
self-consistency condition, the free-energy expansion, the susceptibility and the
phase structure. This stands on its own and does not depend on any experimental
outcome. One coefficient error in the proposal's Equation~(10) is corrected in the
derivation.

\item \textbf{A measured result on 400 items.} One hundred items in each of four
categories, length-matched, with wire-service tells removed and sources balanced
within each category, so that the categories differ by content rather than by where
they came from.

\item \textbf{Two findings that are new to the literature.} First, that the released
TRIBE v2 checkpoint predicts only surface activity, which constrains any study built
on it. Second, that ratio-type measures over standardised model output are undefined
on real content. This is the first independent research application of this model.
\end{enumerate}

\subsection{Not delivered, and why}

\begin{enumerate}[leftmargin=1.6em]
\item \textbf{Any statement about the amygdala.} The instrument cannot reach it.
\item \textbf{A calibrated value for $\hat{\alpha}$.} The fit is not distinguishable
from zero, so quoting a number would misrepresent it.
\item \textbf{1{,}500 items across text, audio and video.} The corpus is 400 text
items. Each item takes about 75 seconds of GPU time, and the available quota does
not cover the original size. Text is the modality the proposal processes first by
design.
\item \textbf{Research Question IV.} It compares information across modalities,
which requires the audio and video that were dropped.
\end{enumerate}

\subsection{Status of each objective}

\begin{center}
\begin{tabular}{@{}clp{7.6cm}@{}}
\toprule
\textbf{\#} & \textbf{Objective} & \textbf{Status after this amendment} \\
\midrule
(i)   & Build the corpus        & Delivered at 400 text items, not 1{,}500 multimodal \\
(ii)  & Run TRIBE, compute the index & Delivered, using the signed index \\
(iii) & Characterise the distribution & Delivered \\
(iv)  & Calibrate $\hat{\alpha}$ & Attempted; reported as a null result \\
(v)   & Derive the phase structure & Delivered across a range of $\alpha$, not at one fitted value \\
(vi)  & Validate as a classifier & Delivered, reporting AUC as the headline figure \\
(vii) & Deliver the instrument  & Delivered; this is the primary deliverable \\
\bottomrule
\end{tabular}
\end{center}

\subsection{What an examiner receives}

A thesis that sets out a physics framework, builds the instrument the proposal
promised, runs it on a controlled corpus, and reports honestly what it found and
what it could not reach. The result may well be that the surface-level brain proxy
does not separate manipulative from neutral content. If so, that is the finding,
reported with the sample size and the smallest effect the design could have detected.

\section{Decisions requested}

\begin{enumerate}[leftmargin=1.6em]
\item Approve the amended measurement (cortical proxy, signed difference) and the
corresponding rewrite of the abstract, \S4.1 and Equation~(4). \emph{Blocking.}
\item Approve reducing the corpus to 400 text items and dropping audio and video,
with Research Question IV declared out of scope. \emph{Blocking.}
\item Confirm that the null calibration is reported as the result, with no value of
$\hat{\alpha}$ quoted anywhere. \emph{Blocking.}
\item Note the substitution of SemEval-2019 Task 4 for the withdrawn NELA-GT-2021.
\end{enumerate}

\vspace{1.5em}
\begin{tabular}{@{}p{7cm}p{7cm}@{}}
\hrulefill & \hrulefill \\
Brian Mwai (1050555) \quad Date & Dr. Songa Mutambi \quad Date \\
\end{tabular}

\end{document}
